\documentclass[book.tex]{subfiles}
\begin{document}

\chapter{The Finite Difference Method}

\label{chap:linear_boundary}
\noindent\rule{11cm}{0.4pt} \\
\textbf{Prerequisites:} \autoref{chap:linear_systems} \\
\textbf{Difficulty Level:} ** \\
\noindent\rule{11cm}{0.4pt}
\vspace{5mm}

In this chapter we learn how to implement the finite-difference method and study  how derivative boundary conditions are incorporated into the finite-difference method.
By the end of the chapter, you will know how to solve 2nd-order nonlinear ODEs with the finite-difference method by using root-location methods for systems of nonlinear algebraic equations.


\section{Finite Difference Method}

A boundary value problem is a differential equation together with a set of additional constraints, called the boundary conditions. A solution to a boundary value problem is a solution to the differential equation which also satisfies the boundary conditions. In the finite difference method, finite differences are substituted for the derivatives in the original equation. Thus, a linear differential equation is transformed into a set of simultaneous algebraic equations which can then be solved. 

Consider a pod of mass $m$ traveling inside a vacuum tube which is 1 mile in length. The pod experiences a drag force proportional to its velocity (by factor $k$) due to the magnets used to levitate and brake. The vacuum tube is maintained at very low pressure (almost close to vacuum) and hence the drag due to air is almost zero. A pusher attached to the pod provides a constant propelling force $F$ to move the pod in the tube. Friction due to wheels is neglected as the pod is assumed to levitate. This setup can be modeled as 
\begin{align}
m\dfrac{d^2x}{dt^2} = F - k\dfrac{dx}{dt}
\end{align}

\begin{marginfigure}[\baselineskip] 
\centering
% \includegraphics[width = 2.4in]{figures/part1b/boundary_value/hyperloop.PNG}
\includegraphics[width = 2.4in]{figures/part1b/boundary_value/Hyperloop_Finite_approach.png}
\centering
\caption{For finite-difference approach, time is divided into a series of time steps.}
\label{fd:fig:1}
\end{marginfigure}

As shown in Figure \ref{fd:fig:1}, the solution domain is first divided into a series of time steps. At each time step, finite difference approximations can be written for the derivatives in the equation. For example, at time step $i$, the second and first derivative can be represented by 
\begin{align}
\frac{d^2x}{dt^2} &= \dfrac{x_{i-1}-2x_i+x_{i+1}}{\Delta t^2}\\
\frac{dx}{dt} &=  \dfrac{x_{i+1}-x_i}{\Delta t};
\end{align}

This approximation can be substituted into the governing equation to give
\begin{align}
m\left (\dfrac{x_{i-1}-2x_i+x_{i+1}}{\Delta t^2}\right) &= F - k\left (\dfrac{x_{i+1}-x_i}{\Delta t}\right )\\
-mx_{i-1} + (k\Delta t + 2m)x_i - (k\Delta t + m)x_{i+1} &= -F\Delta t^2
\end{align}

If the domain is divided into $n$ time steps then $x_0$ and $x_n$ values are specified by the boundary conditions. Therefore, the problem reduces to solving $n-1$ simultaneous linear algebraic equations for the $n-1$ unknowns. The time steps are numbered consecutively, and since each equation
consists of a time step $(i)$ and its adjoining neighbors ($i-1$ and $i+1$), the resulting set of linear algebraic equations will be tridiagonal. \\


\section{Tridiagonal Matrix Algorithm}
This method, also known as the Thomas algorithm, is a simplified form of Gaussian elimination that can be used to solve tridiagonal systems of equations. A tridiagonal system for $n$ unknowns may be written as

\begin{align}
a_ix_{i-1} + b_ix_i + c_ix_{i+1} = d_i
\end{align}
where $a_1 = 0$ and $c_n = 0$. 
\begin{margintable}
\[
\begin{bmatrix}
b_1 & c_1 &  &  & 0\\ 
a_2 & b_2 & c_2 &  & 0 \\ 
  & a_3 & b_3 & \ddots & \\
 &  & \ddots & \ddots & c_{n-1}\\
0  &  &  & a_n & b_n
\end{bmatrix} 
\left \{
\begin{tabular}{c}
$x_1$ \\
$x_2$ \\
$x_3$  \\
$\vdots$ \\
$x_n$
\end{tabular}
\right \}
=
\left \{
\begin{tabular}{c}
$d_1$ \\
$d_2$ \\
$d_3$ \\ 
$\vdots$\\
$d_n$
\end{tabular}
\right \}
\]
\end{margintable}
For such systems, the solution can be obtained in \textit{O(n)} operations instead of \textit{$O(n^3)$} required by Gaussian elimination. A first sweep eliminates the $a_i$'s, and then an (abbreviated) backward substitution produces the solution.\\
The equations to be solved are:
\begin{align}
 b_0 x_0 + c_0x_1 = d_0 \qquad &i = 0\\
a_ix_{i-1} + b_ix_i + c_ix_{i+1} = d_i \qquad &i = 1,....,n-2\\
a_{n-1}x_{n-2} + b_{n-1}x_{n-1} = d_{n-1}\qquad  &i = n - 1
\end{align}
Multiplying the first equation with the fraction ($a_1/b_0$) and subtracting it from the second equation $(i = 1)$ to eliminate $x_0$ gives 
\begin{align}
\left (b_1 - \frac{a_1}{b_0}c_0\right )x_1 + c_1x_2 = d_1 - \frac{a_1}{b_0}d_0
\end{align}


Next modify the third equation $(i = 2)$ with the initial second equation to eliminate $x_1$. This procedure is repeated until the $(N-1)^{th}$ row; the (modified) $(N-1)^{th}$ equation will involve only one unknown, $x_{N-1}$. This may be solved for and then used to solve the $(N-2)^{th}$ equation, and so on until all of the unknowns are solved for. The coefficients of the modified equations can be defined as follows:
\begin{align}
b_i^{'} &= b_i - \left (\dfrac{a_i}{b_{i-1}}\right )c_{i-1}\\
c_i^{'} &= c_i\\
d_i^{'} &= d_i - \left (\dfrac{a_i}{b_{i-1}}\right )d_{i-1}
\end{align}

where the system is redefined as:
\begin{align}
b'_ix_i + c_ix_{i+1} &= d'_i \ \mathrm{for}\ i = 1,\dotsc,N-2\\
b'_{N-1}x_{N-1} &= d'_{N-1} \ \mathrm{for}\  i = N-1
\end{align}
The last equation involves only one unknown. Solving it in turn reduces the next last equation to one unknown, so that this backward substitution can be used to find all of the unknowns:
\begin{align}
b'_ix_i + c_ix_{i+1} &= d'_i \  \mathrm{for}\  i = n-1,n-2,\dotsc,1
\end{align}
Let us now consider how to implement the Tridiagonal matrix algorithm in Physika


%m = 150 #mass in kg
%k = 5 #Magnetic Drag coefficient in Ns/m 
%F = 500 #Force from pusher in N
%  x_i = 0
%  x_e = 490
\begin{lstlisting}[language=python]
def tridiagonal(F: $\mathbb{R}$, m: $\mathbb{R}$, k: $\mathbb{R}$, (xi, xe) : $\mathbb{R}^2$, (ti,te): $\mathbb{R}^2$, dt: $\mathbb{R}$):
  N = (te - ti)/dt- 1
  a = -m*ones(N)
  b = (k*dt + 2*m)*ones(N)
  c = -(k*dt + m)*ones(N)
  d = (-F*dt$^2$)*ones(N)
  d[0] = -F*dt$^2$ + m*xi
  d[N-1] = -F*dt$^2$ + xe*(k*dt + m)
  x = ones(N)

  for i:
    b[i] = b[i] - (a[i]/b[i-1])*c[i-1]
    d[i] = d[i] - (a[i]/b[i-1])*d[i-1]  
  x[N-1] = d[N-1]/b[N-1] # Perform the backsubstitution
  for j = N-2:-1:0:
    x[j] = (d[j] - x[j+1]*c[j])/b[j]
  x
\end{lstlisting}

We have implemented this equation for the specific equation 
\begin{align}
    m\frac{d^2x}{dt^2} &= f - k \frac{dx}{dt}.
\end{align}
This implementation will not directly generalize to higher order equations, which will require a more complex equation solution method.

\section{Derivative Boundary Conditions}
There are three types of boundary conditions commonly encountered in the solution of ordinary or partial differential equations:

\begin{itemize}
	\item \textit{Dirichlet boundary condition}: It specifies the values that a solution needs to take on along the boundary of the domain.
	\item \textit{Neumann boundary condition}: It specifies the values that the derivative of a solution is to take on the boundary of the domain.
	\item \textit{Robin boundary condition}: It is a specification of a linear combination of the values of a function and the values of its derivative on the boundary of the domain.
\end{itemize}
The derivation in the previous section can be modified to work for Neumann boundary conditions.
%Consider the previous example to demonstrate how a derivative boundary condition can be incorporated into the finite-difference approach. In this case assume that the pusher is detached and hence $F = 0$ N.

%\begin{align}
%m\dfrac{d^2x}{dt^2} &=  k\dfrac{dx}{dt}
%\end{align}
%We will prescribe a Neumann boundary condition at one end of the time domain:

%\begin{align}
%\dfrac{dx}{dt}(15) = 42.823 \mathrm{ms}^{-1}, x(30) = 1015
%\end{align}

%Therefore the boundary conditions are applied to obtain linear equations. Now, 

%\begin{align}
%\dfrac{dx}{dt}(15) &= \dfrac{x_1 - x_0}{\Delta t}\\
%x_0 &= x_1 - \Delta t \dfrac{dx}{dt}
%\end{align}

%Therefore on substituting the value of $x_{0}$ in the equation at time step $i = 1$ we get $ 165 x_1 - 165x_2  = -19270.35$

%Consequently, we have incorporated the derivative into the balance.\\

%\textit{Example}: Solve the above problem (with Neumann boundary condition) to obtain distance traveled at the end of each time step.\\
%%\textit{Solution}: 
%\begin{align}
%\begin{bmatrix} 
%165 & -165 & 0 & 0\\ 
%-150 & 315 & -165 & 0\\
%0 & -150 & 315 & -165\\
%0 & 0 & -150 & 315\\
%\end{bmatrix} 
%\left \{
%\begin{tabular}{c}
%$x_1$ \\
%$x_2$ \\
%$x_3$ \\
%$x_4$
%\end{tabular}
%\right \}
%=
%\left \{
%\begin{tabular}{c}
%-19270.35 \\
%0 \\ 
%0 \\
%167475
%\end{tabular}
%\right \}
%\end{align}
%These equations can be solved for
%\begin{align}
%\begin{gathered}
%x_1 = 607.7706,
%x_2 = 724.5606,
%x_3 = 830.7333,
%x_4 = 927.2539
%\end{gathered}
%\end{align}
%Therefore 
%\begin{align}
%\begin{gathered}\
%x_0 = 607.7706 - (3)(42.823) = 479.3016
%\end{gathered}
%\end{align}
%Analytical solution:
%\begin{align}
%x(t) = (-2225.77) e^{-t/30} + (1833.82) 
%\end{align}
%\begin{marginfigure}
%	\begin{center}
%		\begin{tabular}{ | p{0.5cm} | p{2.5cm} | p{2cm} |}
%			\hline
%			\textbf{$t$(s)} & \textbf{$x$(m) Anal. Soln} & \textbf{$x$(m) FDM} \\ \hline
%			15 & 483.3223 & 479.3016 \\ \hline
%			18 & 611.7915 & 607.7706 \\ \hline
%			21 & 728.0353 & 724.5606\\ \hline
%			25 & 833.2171 & 830.7333 \\ \hline
%			27 & 928.3894 & 927.2539 \\ \hline
%			30 & 1015 & 1015 \\ \hline
%		\end{tabular}
%	\end{center}
%\end{marginfigure}
%Table provides a comparison between the analytical solution and the numerical solution obtained with the finite difference method.
%\end{document}
\section{The Finite Difference Method for Nonlinear ODEs}

For nonlinear ODEs, the substitution of finite differences yields a system of nonlinear simultaneous equations. Thus, the most general approach to solving such problems is to use
root-location methods for systems of equations such as the Newton-Raphson method which will be introduced in next lecture. An adaptation of successive substitution can sometimes provide a simpler alternative.

Consider the tube problem solved above. If the pressure in the tube is not close to vacuum then the air drag comes into picture. 
\begin{align}
m\dfrac{d^2x}{dt^2} = F - k\dfrac{dx}{dt} - \dfrac{1}{2}\left (\dfrac{dx}{dt}\right )^2C_DA
\end{align}
We can convert this differential equation into algebraic form by writing it for a time step $i$ and substitute for the second derivative.
\begin{align}
m\left (\dfrac{x_{i-1}-2x_i+x_{i+1}}{\Delta t^2}\right) = F - k\left (\dfrac{x_{i+1}-x_i}{\Delta t}\right )-\frac{1}{2}\left (\dfrac{x_{i+1}-x_i}{\Delta t}\right)^2C_DA
\end{align}
 For routinely encountered problems in engineering and science, if we assume that the unknown nonlinear term is equal to its value from the previous iteration, then the equation can be solved and iterate until the process converges to an acceptable tolerance. Although this approach will not work for all cases, it converges for many ODEs derived from physically based systems.
%\end{document}
%\section{Exercises}
%\begin{enumerate}
%    \item \textcolor{red}{TODO}
%\end{enumerate}
\end{document}